\section{Introduction}
An Agent receives the instruction \emph{process the next support ticket}. It reads a customer's request, checks an order and approval, consults a model, and submits a refund. Encryption hides messages from a network observer but leaves contacted services, traffic sizes and timing visible. The tool's answer also becomes input to another service: the model. Adding a decoy request has a second consequence if the Agent executes it: the defense now changes the work done for the user. Figure~\ref{fig:chain} follows the ticket through these services.

\begin{figure*}[t]
\centering
\begin{tikzpicture}[x=1cm,y=1cm,font=\small,>=Latex]
\tikzset{box/.style={draw,rounded corners,align=center,minimum height=.85cm,text width=2.4cm},
 note/.style={align=center,font=\footnotesize,text width=3.1cm}}
\node[box,fill=blue!5] (u) at (1.3,1.15) {User\\``Process the next ticket''};
\node[box,fill=blue!5] (a) at (5.1,1.15) {Agent\\plan, check, act, retry};
\node[box,fill=green!6] (m) at (9.2,2.6) {Model service\\next action / streamed output};
\node[box,fill=green!6] (t) at (9.2,-.35) {Tool service\\ticket, order, refund};
\node[box,fill=orange!8] (d) at (14.1,1.15) {Specialist Agent\\optional delegation};
\draw[<->,thick] (u)--node[above,yshift=15pt]{client link C}(a);
\draw[<->,thick] (a.north east)--node[above,sloped]{model link M}(m.west);
\draw[<->,thick] (a.south east)--node[below,sloped]{tool link T}(t.west);
\draw[<->,thick] (a.east)--node[above]{coordination A}(d.west);
\node[note] at (2.0,-.35) {\textbf{Transit observer}\\lengths, gaps, connections;\\destinations when available};
\node[note] at (13.7,-.35) {\textbf{Recipient / active service}\\received arguments;\\optionally selected responses};
\node[draw=black!45,fill=black!3,align=left,text width=15.7cm,inner sep=6pt] at (7.9,-1.65) {
\textbf{Two protection questions:} Where does tool feedback appear next? What work does privacy cover execute before it is discarded? The original task and each observable service remain part of the protection decision.};
\end{tikzpicture}
\caption{Running support-ticket example and observation rights. A tool result returns to the Agent and can enter the next model request. Each encrypted link exposes its own traffic shape; a recipient additionally sees the request addressed to it. Concurrent tickets may share a connection on either link. The arrows represent application communication, not a one-to-one mapping from tool events to packets.}
\label{fig:chain}
\end{figure*}


\textbf{The unit of protection is the execution across services.} Classical traffic analysis provides observer models, segmentation methods and adaptive evaluation. We ask which conditions these results need when applied to Agent execution. Twenty-two mechanism anchors and fifteen direct Agent studies supply two complementary layers: the established relation and the evidence for its Agent application. Forty-four setting records connect them to the interfaces in Figure~\ref{fig:chain}.

\paragraph{Feedback connects separately protected services}
Research queries can be recovered from tool destinations, and task labels can be recognized on model connections~\cite{jeong2026wra,li2026ant,pouryousef2026traffic}. A contacted tool has additional access to its arguments. These observations put protections at different boundaries: tool-response padding leaves the model's later request outside its scope, while transit encryption leaves the recipient's arguments visible. Our map links each inference to its access rights and execution stage. The resulting design question is which representations remain exposed after a defense, including initial instructions, forwarded feedback and subsequent actions.

\paragraph{Executed cover needs an application constraint}
Cover discarded at a transport endpoint consumes bandwidth; a decoy executed by an Agent also consumes application resources and may enter later decisions. In the research-Agent study WRA, five decoys perform additional research and lower query recovery, but the reported task-utility scores belong to a separate visibility intervention~\cite{jeong2026wra}. In the live Oblivious HTTP (OHTTP) study, adapted inference is paired with bytes and delay~\cite{dangol2026workflow}. Reading these settings together identifies the missing measurement: the original task's outcome under the same intervention that produces the privacy result. This distinction supports a design space organized by where cover terminates and which application effects it permits.

Website fingerprinting (WF) and machine-learning security establish realistic populations, segmentation, transfer and adaptive evaluation~\cite{juarez2014critical,wang2016realistic,dyer2012peekaboo,mathews2023critical,arp2022dos,tramer2020adaptive}. We organize their application around three contributions. A \emph{transfer matrix} connects each established condition to the Agent setting that preserves it, adds an obligation or leaves it untested. The \emph{deployment and defense maps} expose the two execution boundaries above and the rights needed to observe them. An \emph{evidence-availability summary} uses question-specific denominators: four query/workflow privacy settings contain two located adaptation results and one matched original-task utility result, with no setting containing both. Identity obfuscation and adversarial-Agent blocking retain their separate objectives.

A paired mechanism study complements the synthesis with identical initial messages, category-bearing tool feedback and randomized task composition over LangGraph, Model Context Protocol (MCP) Streamable HTTP and streaming HTTP/2. Its 3,120 local executions and 192 hosted-model attempts separate first-feedback recognition, complete-execution exposure and task discovery. All twelve adjusted primary comparisons exceed .05, including post-feedback versus prefix. The pilot supplies bounded measurements; the transfer matrix states which design conclusions follow from the source evidence and which need a further test.
